\chapter{Arithmetic \(E_{n}\) Formality and CoHA}
\label{v4-arith:kontsevich-soibelman}

\section{Rectification of the ordered product}
\label{v4-arith:ks:E1-automatic}

The little-interval operad controls an ordered multiplication before
passing to cohomology. Its integral formality permits associative DGA
models. Formality of a particular DGA asks for a further multiplicative
comparison with its cohomology.

\begin{theorem}[integral formality of the little-interval operad]
\label{v4-arith:thm:E1-formality-automatic}
\ClaimStatusProvedHere{}
For every commutative ring $S$, the component augmentation gives an
operadic quasi-isomorphism
\[
 C_\bullet(E_1;S)\longrightarrow\mathrm{Ass}_S.
\]
Here $\mathrm{Ass}_S(n)=S[\Sigma_n]$ records the order of $n$ inputs.
\end{theorem}

\begin{proof}
Each component of the space of $n$ disjoint labeled intervals is
specified by their left-to-right order.
For a fixed order, their endpoints satisfy linear inequalities.
Their parameter space is convex and therefore contractible.
There are $n!$ components, freely permuted by $\Sigma_n$.
The singular-chain augmentation maps every zero-simplex to its
component and every positive-dimensional simplex to zero.
It is a quasi-isomorphism in each arity.
Substitution of intervals substitutes their input orders, so these
maps respect operadic composition and units.
\end{proof}

\begin{corollary}[rectification and cohomology]
\label{v4-arith:cor:ordered-E1-survives}
\ClaimStatusProvedHereConditional{}
Whenever the operadic change-of-algebra adjunction satisfies its
rectification hypotheses, an ordered $E_1$-algebra has an associative DGA model.
This conclusion concerns its multiplication on chains.
It supplies no quasi-isomorphism to its cohomology algebra.
The nonformal algebras of Theorem~\ref{v4-arith:ks:coefficient-formality}
are strict associative examples over $\mathbb Z$.
\end{corollary}

\begin{proof}
Apply the rectification adjunction to the preceding operadic quasi-isomorphism.
A DGA still has its differential, and passage to cohomology requires a
separate multiplicative quasi-isomorphism zigzag.
For the last assertion take $S=\mathbb Z$, $c=2$, and $r=2$ in
Theorem~\ref{v4-arith:ks:coefficient-formality}.
\end{proof}

\section{An integral obstruction for the little-disk operad}
\label{v4-arith:ks:E2-fails}

The action of input permutations is part of a symmetric operad.
It already distinguishes chains on two little disks from their homology.

\begin{theorem}[the arity-two obstruction]
\label{v4-arith:thm:E2-fails-integrally}
\ClaimStatusProvedHere{}
The chain operad $C_\bullet(E_2;S)$ is not formal as a symmetric
operad for $S=\mathbb Z$ or $S=\mathbb F_2$.
\end{theorem}

\begin{proof}
In arity two, the direction between the disk centres gives an
equivariant homotopy equivalence with $S^1$.
The transposition group $G=\Sigma_2$ acts by the antipodal map.
This action is free and induces the identity on both homology groups.
Thus homology with zero differential is
$S\oplus S[1]$ as a complex of $S[G]$-modules, in homological grading.

Derived coinvariants of the actual chains compute
$H_\bullet(S^1/G;S)$, since the action is free.
They vanish above degree one.
Derived coinvariants of the zero-differential homology complex are
\[
 H_\bullet(G;S)\oplus H_\bullet(G;S)[1].
\]
The periodic resolution for $G=\langle g:g^2=1\rangle$ has
alternating maps $1-g$ and $1+g$.
On the trivial module they become $0$ and $2$.
For $S=\mathbb Z$, this gives a nonzero $\mathbb Z/2$ in every positive odd degree.
For $S=\mathbb F_2$, it gives a nonzero group in every degree.
Hence the two derived coinvariant complexes cannot be quasi-isomorphic.
An operadic formality zigzag would give such an equivalence already in arity two,
which is impossible.
\end{proof}

\begin{corollary}[coefficient and operad choices]
\label{v4-arith:cor:no-integral-E2}
\ClaimStatusProvedHere{}
The integral chain operad cannot be replaced by
$\mathrm{Ger}_{\mathbb Z}$ through a symmetric-operad formality zigzag.
Integral $E_2$-algebras can nevertheless exist.
Over $\mathbb Q$, Tamarkin's formality theorem supplies a formality
comparison for the operad, with a choice of such comparison
(Tamarkin, \emph{Formality of chain operad of little discs},
Lett.~Math.~Phys.~66 (2003), Theorem~2.1, p.~67).
This operadic comparison does not make every individual algebra formal.
\end{corollary}

\begin{proof}
The integral assertion follows from the preceding theorem, since
$H_\bullet(E_2;\mathbb Z)=\mathrm{Ger}_{\mathbb Z}$.
An algebra over the chain operad requires an operad action, which does
not require formality of the acting operad.
For example, an ordinary commutative algebra is an algebra over that
operad through its augmentation to the commutative operad.
The rational assertion is the cited operadic theorem.
Its change-of-operad construction retains the differential of each algebra.
\end{proof}

\begin{remark}[obstructions under coefficient change]
The class of a higher multiplication belongs to a specified Hochschild
cochain complex. Comparing such classes after coefficient change requires
a comparison of those complexes.
For the explicit family below, monic division supplies this comparison
on the chosen minimal structure.
Its obstruction is detected directly in $S/cS$ by
Theorem~\ref{v4-arith:ks:coefficient-formality}.
No quotient of unrelated Hochschild cohomology groups is used.
\end{remark}

\section{Strict operations and transferred operations}
\label{v4-arith:ks:higher-products-vanish}

Let $D_{\mathrm{ord}}$ be an associative DGA representing an ordered
arithmetic factorization algebra, when that representation has been constructed.
Its differential and product are its operations on chains.
The Hochschild trace class $\mathcal V^{\mathrm{prim}}$ remains a separate
object until a comparison with that algebra is specified.

\begin{theorem}[operations in a strict model]
\label{v4-arith:thm:mk-vanish}
\ClaimStatusProvedHere{}
Every associative DGA $D_{\mathrm{ord}}$ defines an $A_\infty$-algebra
with $m_1=d$, $m_2$ its product, and $m_k=0$ for $k\geq3$.
Its transferred operations on cohomology need not vanish.
\end{theorem}

\begin{proof}
The arity-one identity is $d^2=0$.
The arity-two identity is the Leibniz rule, and the arity-three
identity is associativity. Every higher identity has zero terms.
The transfer calculation in Lemma~\ref{v4-arith:ks:transfer-four} gives
$m_4=-uF\smile F$ on cohomology of a strict DGA.
Its nonzero obstruction class for a nonunit parameter is proved in
Theorem~\ref{v4-arith:ks:coefficient-formality}.
\end{proof}

\begin{corollary}[the two formality questions]
\label{v4-arith:cor:strict-associativity-Z}
\ClaimStatusProvedHere{}
A strict integral model can have a nonformal minimal model on its cohomology.
Therefore the vanishing of its displayed chain operations for $k\geq3$
does not remove higher-operation obstructions to arithmetic centre comparisons.
\end{corollary}

\begin{proof}
The model $A(\mathbb Z,2,2)$ has precisely these properties by
Theorem~\ref{v4-arith:ks:coefficient-formality}.
\end{proof}

\section{A relative derived centre and its nonpositive part}
\label{v4-arith:ks:relative-centre}

A strict associative differential graded algebra can retain higher
multiplications on its cohomology. Take $\mathbb Z[t]\langle\tau\rangle$, with central $t$ and
\[
 |t|=0,\qquad |\tau|=-1,\qquad d\tau=t^2.
\]
The relation $[t]^2=0$ has primitive $\tau$, whose square
$u=\tau^2$ is a nonzero class of degree $-2$.
The fourth operation will record this square.
The same calculation occurs inside a relative derived centre, where
an additional class of degree two changes the multiplication relation.

Fix a commutative unital ring $S$ and an integer $r\geq1$.
All differentials have cohomological degree one.
Algebras are unital and associative, with a central closed $S$-action.
Tensor products and polynomial sums are algebraic.
A homogeneous tensor of maps acts by
\[
 (f\otimes g)(x\otimes y)=(-1)^{|g||x|}f(x)\otimes g(y).
\]
An odd generator need not have square zero in this category.
Formality over $S$ means a finite zigzag of unital $S$-linear
DGA quasi-isomorphisms to cohomology with zero differential.
No additional generator is fixed by this definition.

Put $R=S[t]$ and
\[
 E_r=R\langle\tau\rangle,\qquad d\tau=t^r.
\]
The variable $t$ commutes with $\tau$.
Over the characteristic-zero geometric base, this is the algebra in
Booth, \emph{The derived contraction algebra}, Lemma~4.3.8, p.~22,
with $(\xi,\zeta,n)=(t,\tau,r)$.
The following calculation uses only its displayed algebraic presentation.

The relative derived centre is
\[
 Z_R(E_r)=\operatorname{RHom}_{E_r\otimes_R E_r^{\mathrm{op}}}
 (E_r,E_r),
\]
with its derived composition product.
The opposite algebra uses the Koszul sign convention.
The word ``relative'' fixes $R=S[t]$ in the bimodule tensor product.
In particular, this construction differs from $Z_S(E_r)$.

\begin{proposition}[a multiplicative model for the relative centre]
\label{v4-arith:ks:centre-model}
\ClaimStatusProvedHere{}
The relative centre has the associative DGA model
\[
 C_{S,r}=S[t]\langle\tau\rangle[z]/(z^2),\qquad
 |z|=2,\qquad
 D\tau=t^r-2\tau^2z,\quad Dt=Dz=0,
\]
where $z$ commutes with $t,\tau$.
There is a multiplicative quasi-isomorphism from this algebra to the
endomorphism DGA of a free bimodule resolution of $E_r$.
\end{proposition}

\begin{proof}
Let $P$ be the free graded $E_r$-bimodule on generators
$e_0,e_1$ of degrees $0,-2$, with
\[
 d_Pe_0=0,\qquad d_Pe_1=\tau e_0-e_0\tau.
\]
The Leibniz rule defines the differential on all of $P$.
Its square on $e_1$ is $t^re_0-e_0t^r=0$.
Define $\epsilon:P\to E_r$ by
$\epsilon(e_0)=1$ and $\epsilon(e_1)=0$.

The tensor-algebra sequence
\[
 0\longrightarrow E_r\otimes_R R\tau\otimes_R E_r
 \longrightarrow E_r\otimes_R E_r
 \xrightarrow{\mathrm{mult}} E_r\longrightarrow0
\]
is exact as a sequence of graded modules.
The first map sends $a\otimes\tau\otimes b$ to
$a\tau\otimes b-a\otimes\tau b$.
At each total word length, its matrix is the difference map between
successive cuts of one word.
Its image consists exactly of combinations whose coefficients sum to zero,
and its kernel vanishes by inspecting the first cut, then each successive cut.
The universal differential of $d\tau=t^r$ relative to $R$ is zero.
Thus this sequence is a sequence of complexes.
Its two-term totalization, with the usual shift signs, is $P$.
Exactness proves that $\epsilon$ is a quasi-isomorphism.
The filtration generated first by $e_0$ and then by $e_1$ makes $P$
a semifree bimodule, so it computes the displayed derived Hom.

There is a counital coassociative diagonal
$\Delta:P\to P\otimes_{E_r}P$ given by
\[
 \Delta(e_0)=e_0\otimes e_0,\qquad
 \Delta(e_1)=e_1\otimes e_0+e_0\otimes e_1.
\]
The two middle terms in $d\Delta(e_1)$ cancel in the balanced tensor product.
This proves that $\Delta$ commutes with the differential.
Coassociativity and the counit identities follow on these two generators.

A homogeneous bimodule map of degree $k$ satisfies
$f(apb)=(-1)^{k|a|}af(p)b$.
Identify $\operatorname{Hom}_{E_r^e}(P,E_r)$ with $E_r\oplus E_rz$
by the values on $e_0,e_1$.
Convolution through $\Delta$ gives
\[
 (a+bz)(a'+b'z)=aa'+(ba'+ab')z.
\]
The differential on the value $a$ has additional $e_1$-component
\[
 -\tau a+(-1)^{|a|}a\tau.
\]
For $a=\tau$, this is $-2\tau^2$.
These are exactly the multiplication and differential of $C_{S,r}$.
In particular $D(\tau^2)=0$, so $D^2=0$.

For a convolution cochain $f$, the map
\[
 \rho(f)=(1\otimes f)\Delta:P\longrightarrow P
\]
uses the right action to identify $P\otimes_{E_r}E_r$ with $P$.
Coassociativity gives $\rho(fg)=\rho(f)\rho(g)$, with the tensor signs above.
The counit gives $\epsilon\rho(f)=f$.
Since $P$ is semifree, composition with $\epsilon$ is a
quasi-isomorphism from $\operatorname{End}_{E_r^e}(P)$ to
$\operatorname{Hom}_{E_r^e}(P,E_r)$ as complexes.
Its section $\rho$ is therefore the claimed DGA quasi-isomorphism.
\end{proof}

To separate the coefficient in the differential from the base ring,
let $c\in S$ and define
\begin{equation}\label{v4-arith:ks:coefficient-family}
 \begin{gathered}
 B=S[u,v]/(v^2),\quad T=B[t],\quad f=t^r-cv,\\
 A(S,c,r)=T\oplus\tau T,\qquad
 |u|=-2,\quad |v|=|t|=0,\quad |\tau|=-1,\\
 \tau^2=u,\qquad d(a+\tau b)=fb.
 \end{gathered}
\end{equation}
Here $\tau$ commutes with $T$.
The equality $d(\tau^2-u)=f\tau-\tau f=0$ verifies the defining relation.
This family allows arbitrary $c$.
The centre in Proposition~\ref{v4-arith:ks:centre-model} has the specific coefficient $c=2$.

\begin{proposition}[the truncation map and its quotient]
\label{v4-arith:ks:centre-truncation}
\ClaimStatusProvedHere{}
There is an injective DGA map
\[
 j:A(S,2,r)\longrightarrow C_{S,r},\qquad
 j(t)=t,\quad j(\tau)=\tau,\quad j(u)=\tau^2,\quad j(v)=\tau^2z.
\]
Its image is the nonpositive truncation $\tau_{\leq0}C_{S,r}$.
It induces an isomorphism on cohomology in every nonpositive degree.
The quotient complex has its only cohomology in degree two, where it is
$S[t]/(t^r)\,z$.
\end{proposition}

\begin{proof}
Write $u=\tau^2$ inside $C_{S,r}$.
Its even and odd parts are
\[
 S[u,z,t]/(z^2)\quad\hbox{and}\quad
 \tau S[u,z,t]/(z^2).
\]
The monomials $u^at^i$ and $u^azt^i$ are $S$-linearly independent.
Sending $u^avt^i$ to $u^{a+1}zt^i$ proves injectivity of $j$,
also on the odd summand.
The map respects the differential because $t^r-2v$ maps to $D\tau$.
All degree-zero elements of $C_{S,r}$ are closed.
Its only positive terms are $S[t]\tau z$ in degree one and $S[t]z$ in degree two.
The quotient of $j$ is therefore the complex
\[
 S[t]\tau z\xrightarrow{\ t^r\ }S[t]z
\]
in degrees one and two.
Multiplication by $t^r$ is injective, which proves the assertions.
\end{proof}

\begin{remark}[the missing positive generator]
For $S\ne0$, the map $j$ is not a quasi-isomorphism.
It has no graded algebra retraction: every such retraction sends $z$ to zero,
but a retraction would also send $uz$ to the nonzero element $v$.
The multiplicative quasi-isomorphism of Proposition~\ref{v4-arith:ks:centre-model}
identifies the full $C_{S,r}$ with the relative centre.
The map $j$ identifies $A(S,2,r)$ with its nonpositive truncation.
\end{remark}

\section{The fourth operation over a coefficient ring}
\label{v4-arith:ks:fourth-operation}

The relation $t^r=cv$ has a chosen primitive $\tau$.
Products of two such primitives produce $u$.
To decide whether changing the primitives can remove this effect,
we first compute the transferred operation and then its cohomology class.

Multiplication by the monic polynomial $f=t^r-cv$ on $T$ is injective.
A nonzero polynomial times $f$ retains its nonzero leading coefficient,
including when $S$ has zero divisors.
It follows that
\begin{equation}\label{v4-arith:ks:family-cohomology}
 H=H^*(A(S,c,r))=S[u,v,t]/(v^2,t^r-cv),\qquad H^{\mathrm{odd}}=0.
\end{equation}
Division by the monic polynomial supplies the $S$-basis
\[
 u^av^jt^i,\qquad a\geq0,\quad j\in\{0,1\},\quad 0\leq i<r.
\]
Each $H^{-2a}$ is free of rank $2r$.

Let $i:H\to T$ choose the representative of $t$-degree less than $r$.
Let $p:A\to H$ take the even component modulo $f$ and kill the odd component.
Set
\[
 h(x)=\tau\frac{x-ip(x)}f\quad(x\in T),\qquad h(\tau T)=0.
\]
Existence and uniqueness of monic division give
\[
 dh+hd=1-ip,\qquad pi=1,\qquad h^2=hi=ph=0.
\]
Define the $B$-bilinear degree-zero map $F:H^{\otimes_S2}\to H$ by
\[
 i(x)i(y)-i(xy)=f\,iF(x,y).
\]
On the normal basis it is
\[
 F(t^i,t^j)=
 \begin{cases}
 t^{i+j-r},&i+j\geq r,\\
 0,&i+j<r.
 \end{cases}
\]
For the first case, use
$t^{i+j}=f t^{i+j-r}+cv t^{i+j-r}$.
Both quotients have degree less than $r$.

We use $A_\infty$ operations of degree $|m_n|=2-n$
and morphism components of degree $|f_n|=1-n$.
Explicitly, the operation identities have signs
\[
 \sum_{a+b+q=n}(-1)^{a+bq}
 m_{a+1+q}(1^{\otimes a}\otimes m_b\otimes1^{\otimes q})=0.
\]
The indices satisfy $a,q\geq0$ and $b\geq1$.
In a morphism identity, the term
$m_k(f_{i_1}\otimes\cdots\otimes f_{i_k})$
has sign $(-1)^{\sum_{j=1}^k(k-j)(i_j-1)}$.
Tensor evaluation also uses the Koszul rule fixed above.
A minimal structure has $m_1=0$.

\begin{lemma}[transfer and the first obstruction]
\label{v4-arith:ks:transfer-four}
\ClaimStatusProvedHere{}
The contraction gives a minimal structure on $H$ and a morphism to $A$
whose first components are
\[
 f_1=i,\qquad f_2=-\tau iF,\qquad f_3=m_3=0,
\]
\begin{equation}\label{v4-arith:ks:m-four}
 m_4(x,y,z,w)=-uF(x,y)F(z,w),\qquad
 f_4(x,y,z,w)=\tau u\,iF(F(x,y),F(z,w)).
\end{equation}
The operation $m_4$ is a Hochschild cocycle of bidegree $(4,-2)$.
\end{lemma}

\begin{proof}
For completeness, the transfer can be constructed over $S$ without a field hypothesis.
Set
\[
 U_n=\sum_{a+b=n}(-1)^{a-1}\mu(f_a\otimes f_b),\qquad
 m_n=pU_n,\qquad f_n=-hU_n\quad(n\geq2),
\]
where $a,b\geq1$ and $\mu$ is multiplication in $A$.
Each formula uses only smaller components.

Here is the compatibility argument.
Suspend by $|sx|=|x|-1$, and use the tensor coalgebra
$T^c(sH)=\bigoplus_{n\geq1}(sH)^{\otimes_S n}$ with deconcatenation.
Write $I,P,K$ for the suspended maps $i,p,h$.
Write $D_1$ for the suspended differential of $A$ and $\Lambda$
for its binary multiplication coderivation, with the signs just specified.
Thus $D_1^2=\Lambda^2=D_1\Lambda+\Lambda D_1=0$.
The recursion becomes
\[
 Q_n=\sum_{a+b=n}\Lambda_2(\Phi_a\otimes\Phi_b),\quad
 \Phi_1=I,\quad \Phi_n=-KQ_n,\quad \beta_1=0,\quad\beta_n=PQ_n.
\]
Extend $\Phi$ by consecutive partitions and $\beta$ by insertion.
These give a coalgebra map and a coderivation respectively.

Let $E=(D_1+\Lambda)\Phi-\Phi\beta$.
If $E$ vanishes below length $n$, the coproduct identity
\[
 \Delta E=(E\otimes\Phi+\Phi\otimes E)\Delta
\]
makes $E$ primitive in length $n$.
A primitive element has length one, since the cut after its first factor
detects every component of greater length.
Writing $Q$ for the corestriction of $\Lambda\Phi$, one obtains
\[
 0=\pi\Lambda E=-D_1Q-Q\beta,
 \qquad
 \pi E=K(D_1Q+Q\beta)=0.
\]
The second identity uses $D_1K+KD_1=1-IP$.
Induction gives $E=0$.
The highest-length component of $\Phi$ is $I^{\otimes n}$,
which has left inverse $P^{\otimes n}$.
Thus $\Phi$ is injective on finite tensor sums, and
$\Phi\beta^2=(D_1+\Lambda)^2\Phi=0$ implies $\beta^2=0$.
This proves the transfer identities in every arity.

The maps $p,h$ kill odd outputs, giving $f_3=m_3=0$.
Only the partition $(2,2)$ contributes to $m_4$, with a negative sign.
Its product is $-u\,iF(x,y)iF(z,w)$.
Applying $p$ and $-h$ gives \eqref{v4-arith:ks:m-four}.
For an even algebra the Hochschild differential is
\[
 \begin{split}
 (\delta K)(x_1,\ldots,x_{n+1})={}&x_1K(x_2,\ldots,x_{n+1})\\
 &+\sum_{j=1}^n(-1)^jK(x_1,\ldots,x_jx_{j+1},\ldots,x_{n+1})\\
 &+(-1)^{n+1}K(x_1,\ldots,x_n)x_{n+1}.
 \end{split}
\]
The arity-five identity gives $\delta m_4=0$.
Alternatively, expand the representative of a triple product in both
parenthesizations, cancel the injective factor $f$, and apply $p$.
This gives $\delta F=0$, hence $\delta(F\smile F)=0$.
\end{proof}

A nonzero value of $m_4$ alone does not prove nonformality.
A change of coordinates can alter it by a Hochschild boundary.
An arbitrary formality zigzag must therefore be compared with the
particular minimal structure just constructed.

\begin{lemma}[finite comparison over an arbitrary ring]
\label{v4-arith:ks:projective-transport}
\ClaimStatusProvedHere{}
Suppose a DGA $D$ has an explicit minimal model $V$ whose homogeneous
components are projective $S$-modules.
That model can be transported through any finite quasi-isomorphism zigzag,
through every specified finite arity.
Consequently, formality of $A(S,c,r)$ implies
\[
 m_4=\delta K
\]
for an $S$-linear cochain $K:H^{\otimes_S3}\to H$ of degree $-2$.
The cochain need not preserve $t,u,v$.
\end{lemma}

\begin{proof}
First take a degreewise surjective DGA quasi-isomorphism $q:P\to D$.
Its kernel $L$ is acyclic.
Suppose a coalgebra morphism from the minimal source has been lifted
through arity $n-1$.
Every degree of $W_n=(sV)^{\otimes_S n}$ is projective:
tensor products of projective modules are projective, and these degrees
are direct sums of such tensor products.
Choose a graded lift of the arity-$n$ component through $q$.

The defect $E=D_P\Phi-\Phi b$ is primitive in arity $n$ because
all smaller defects vanish.
It lies in $sL$ because its image under $q$ vanishes.
The identity $D_PE+Eb=0$ makes it closed:
minimality means that $b$ lowers tensor length, so $Eb$ vanishes here.
Acyclicity supplies the surjection
\[
 d:(sL)^j\longrightarrow Z^{j+1}(sL).
\]
Projectivity of $W_n^j$ lifts this particular defect through $d$.
Subtract the lift from the chosen morphism component.
This removes the defect without changing its image under $q$.
No contraction of the acyclic module $L$ is required.

Every backward arrow admits such a surjective replacement.
Let $J$ be the free $S$-module on $e_0,e_1,\eta$, with
$|e_i|=0$, $|\eta|=1$,
\[
 e_i^2=e_i,\quad e_0e_1=e_1e_0=0,\quad
 e_0\eta=\eta e_1=\eta,\quad
 \eta e_0=e_1\eta=\eta^2=0,
\]
\[
 de_0=-\eta,\qquad de_1=\eta,\qquad d\eta=0.
\]
The unit is $e_0+e_1$.
These are upper triangular matrix multiplication and its differential.
The endpoint maps $\mathrm{ev}_i:J\to S$ select the diagonal coefficients.
The map $k(\eta)=e_1$, $k(e_i)=0$ satisfies
\[
 dk+kd=1-\mathrm{const}\,\mathrm{ev}_0.
\]
On $D\otimes_S J$, use $(-1)^{|a|}a\otimes k(\omega)$ on $a\otimes\omega$.
This gives the same contraction with the tensor differential.

For a quasi-isomorphism $g:D'\to D$, form the pullback DGA
\[
 P(g)=\{(b,\omega)\in D'\times(D\otimes_S J):
             \mathrm{ev}_0\omega=g(b)\}.
\]
Projection $p:P(g)\to D'$ is a homotopy equivalence as a complex,
with section $j(b)=(b,g(b)\otimes1)$ and the displayed contraction.
The other endpoint $q=\mathrm{ev}_1$ is degreewise surjective:
$(0,a\otimes e_1)$ lifts every homogeneous $a\in D$.
Since $qj=g$, it is also a quasi-isomorphism.
Thus each backward arrow can be replaced by the span
$D'\xleftarrow{p}P(g)\xrightarrow{q}D$.
Lift through $q$ and compose with $p$.
Forward arrows require only composition.
This traverses the finite zigzag while keeping the minimal source fixed.

At a strict cohomology target, the linear component is a graded algebra automorphism.
Compose with its inverse to make that component the identity.
Evenness of $H$ forces the second component to vanish.
The arity-four morphism equation now reads $m_4=\delta K$,
with $K$ the third component.
\end{proof}

\begin{theorem}[the coefficient criterion for formality]
\label{v4-arith:ks:coefficient-formality}
\ClaimStatusProvedHere{}
For every commutative unital ring $S$, every $c\in S$, and every $r\geq1$,
\[
 \boxed{\quad A(S,c,r)\text{ is formal over }S
       \quad\Longleftrightarrow\quad r=1\ \text{or}\ c\in S^\times.\quad}
\]
\end{theorem}

\begin{proof}[Necessity]
Assume $r\geq2$ and put $a=t^{r-1}$.
The three adjacent products in $(a,t,a,t)$ are $cv$.
Since $F(a,t)=1$, one has
\[
 m_4(a,t,a,t)=-u.
\]
Evaluation $t,v\mapsto0$ defines a graded algebra map $H\to S[u]$.
Follow it in degree $-2$ by extraction of the coefficient of $u$.
The resulting functional $\lambda:H^{-2}\to S$ satisfies
$\lambda(u)=1$ and $\lambda(tH^{-2})=0$.
For every $S$-linear degree-$-2$ ternary cochain,
\[
 \begin{split}
 (\delta K)(a,t,a,t)={}&aK(t,a,t)+K(a,t,a)t\\
 &+c\bigl(-K(v,a,t)+K(a,v,t)-K(a,t,v)\bigr).
 \end{split}
\]
Applying $\lambda$ kills the exterior terms because $a\in tH^0$.
Thus $\lambda(\delta K(a,t,a,t))\in cS$.
The comparison lemma makes formality imply $-1\in cS$, so $c$ is a unit.
Equivalently, evaluation by $\lambda$ modulo $cS$ kills every boundary
and sends $m_4$ to $-1$ in $S/cS$.
This argument reduces one cochain equation modulo $cS$.
It does not tensor an arbitrary quasi-isomorphism with a nonflat quotient.
\end{proof}

\section{A multiplicative resolution in the formal case}
\label{v4-arith:ks:positive-resolution}

For $r=1$, the inclusion $B\hookrightarrow A(S,c,1)$ is a quasi-isomorphism,
since its cohomology is $B[t]/(t-cv)\cong B$.
For unit $c$ and arbitrary $r$, eliminating $v=c^{-1}t^r$ gives
\[
 H\cong S[t,u]/(t^{2r}).
\]
A resolution must account for the square of the primitive of $t^{2r}$.

Put $R'=S[t,u]$ and $s=t^r$.
The only element of degree $-1$ in $A$ whose differential is $s^2$ is
$\tau(s+cv)$, by injectivity of $d$ on the odd summand.
Its square is
\[
 u(s^2+2csv).
\]
For $S\ne0$, this polynomial is nonzero.
Thus the exterior resolution with generator $h$, $h^2=0$, $dh=s^2$
has no DGA map to $A$ over $R'$.
Successive generators remove these successive multiplication defects.

\begin{proposition}[explicit quasi-isomorphisms]
\label{v4-arith:ks:positive-zigzag}
\ClaimStatusProvedHere{}
Let
\[
 P=R'\langle h_1,h_2,\ldots\rangle,\qquad |h_n|=1-2n,
\]
with central $R'$ and differential
\[
 dh_1=s^2,\qquad dh_n=-\sum_{i+j=n}h_ih_j\quad(n\geq2),
\]
where the summation indices are positive.
Define integers $a_1=1$ and $a_n=-\sum_{i+j=n}a_ia_j$.
There are DGA maps fixing $R'$
\[
 R'/(s^2)\xleftarrow{\ \varepsilon\ }P
       \xrightarrow{\ \Phi\ }A(S,c,r),
\]
where
\[
 \varepsilon(h_n)=0,\qquad
 \Phi(h_n)=a_n\tau u^{n-1}\bigl(s+(2n-1)cv\bigr).
\]
The map $\varepsilon$ is always a quasi-isomorphism.
The cohomology map of $\Phi$ in each degree $-2a$ is
$\operatorname{diag}(I_r,cI_r)$.
It is an isomorphism exactly when $c$ is a unit.
\end{proposition}

\begin{proof}
All $h_n$ are odd.
In $d^2h_n$, the scalar terms cancel as
$-s^2h_{n-1}+h_{n-1}s^2$.
Each remaining ordered triple product occurs twice with opposite signs.
Thus $d^2=0$, including in characteristic two.
For $n\geq1$,
\[
 d\Phi(h_n)=a_nu^{n-1}\bigl(s^2+2(n-1)csv\bigr),
\]
because $(s-cv)(s+(2n-1)cv)=s^2+2(n-1)csv$.
For $i+j=n$, one also has
\[
 \Phi(h_i)\Phi(h_j)=a_ia_j u^{n-1}\bigl(s^2+2(n-1)csv\bigr).
\]
The recurrence for $a_n$ proves that $\Phi$ commutes with the differential.
The augmentation $\varepsilon$ does so directly.

Consider the auxiliary DGA
\[
 K=R'\oplus R'h,\qquad |h|=-1,\qquad h^2=0,\qquad dh=s^2.
\]
The exterior relation is imposed also in characteristic two.
The map $\pi:P\to K$ given by $h_1\mapsto h$ and $h_n\mapsto0$ for $n\geq2$
is a DGA map.
Give $h_n$ weight $n$, $h$ weight one, and $R'$ weight zero.
The differential preserves the increasing weight filtration.

At weight $N\geq1$, words $h_{i_1}\cdots h_{i_\ell}$ correspond to
compositions $i_1+\cdots+i_\ell=N$.
Equivalently, they specify the subset of cuts in $\{1,\ldots,N-1\}$.
The associated graded piece is the exterior module
\[
 R'\otimes_S\Lambda_S(e_1,\ldots,e_{N-1}),
 \qquad d_0=-\Bigl(\sum_{j=1}^{N-1}e_j\Bigr)\wedge(-).
\]
A $q$-cut word with coefficient of degree $d$ has degree $d+q+1-2N$.
Splitting the $k$th factor has sign $(-1)^k$,
which is exactly the sign for insertion after $k-1$ existing cuts.
For $N\geq2$, minus removal of $e_1$ contracts this complex:
\[
 \iota_{e_1}(\omega\wedge x)+\omega\wedge\iota_{e_1}(x)=x,
 \qquad \omega=\sum_{j=1}^{N-1}e_j.
\]
This identity uses only the exterior relations and integral signs.
In weights zero and one, $\operatorname{gr}\pi$ is the identity.

The cone of $\pi$ has acyclic associated graded pieces.
Every cone cycle has a greatest weight.
Subtract a differential whose top component is its primitive, and descend in weight.
Finite support makes this process terminate.
Hence $\pi$ is a quasi-isomorphism, without a convergence hypothesis.
Multiplication by the monic element $s^2=t^{2r}$ is injective on $R'$,
so $K\to R'/(s^2)$ is a quasi-isomorphism.
This proves the assertion for $\varepsilon$.

In degree $-2a$, use the source basis
$u^a(1,t,\ldots,t^{2r-1})$ and the target basis
\[
 u^a(1,t,\ldots,t^{r-1},v,vt,\ldots,vt^{r-1}).
\]
The map fixes $t,u$, and $t^{r+i}$ maps to $cvt^i$.
Its matrix is therefore $\operatorname{diag}(I_r,cI_r)$.
Its kernel is $\operatorname{Ann}_S(c)^r$, and its cokernel is $(S/cS)^r$.
A unit $c$ gives the asserted formality zigzag.
Together with the $r=1$ inclusion, this proves sufficiency in
Theorem~\ref{v4-arith:ks:coefficient-formality}.
\end{proof}

\section{The sixth operation on the full centre}
\label{v4-arith:ks:sixth-operation}

The positive generator $z$ changes the obstruction calculation when $2$ is invertible.
It permits a primitive for $m_4$, but that change produces a sixth operation.
A bar cycle containing $z$ detects the resulting class.

Assume throughout this section that $S\ne0$ and $2\in S^\times$.
Write
\[
 H_C=S[u,t,z]/(t^r-2uz,z^2),\qquad |u|=-2,\quad |t|=0,\quad |z|=2.
\]
Its normal basis consists of $u^at^iz^j$ with $a\geq0$, $0\leq i<r$, and $j=0,1$.
Monic division by $t^r-2uz$ gives the same contraction as before.
The transfer lemma, with coefficient algebra $S[u,z]/(z^2)$, gives
\[
 m_4=-uF\smile F,
\]
\[
 m_6(x_1,\ldots,x_6)=u^2\bigl(
 F_{12}F(F_{34},F_{56})+F(F_{12},F_{34})F_{56}\bigr),
 \qquad F_{ij}=F(x_i,x_j).
\]
Indeed, the only nonzero partitions in arity six are $(2,4)$ and $(4,2)$.
Every odd-arity operation vanishes by degree.
In particular the augmentation
$\epsilon:H_C\to S$, $t,u,z\mapsto0$, kills $m_4$ and $m_6$.

To construct a primitive for $m_4$, consider the associative algebra
\[
 S[u,t,z,\alpha,\beta]/(t^r-2uz-\alpha,z^2-\beta),
 \qquad |\alpha|=0,\quad |\beta|=4.
\]
Monic division makes it free over $S[u,\alpha,\beta]$ on $t^iz^j$,
with $0\leq i<r$ and $j=0,1$.
Let $F,G$ be the coefficients of $\alpha,\beta$ in its multiplication,
evaluated at $\alpha=\beta=0$.
They have degrees $0,-4$, respectively.
The first $F$ agrees with the multiplication defect already used.
Associativity in the two parameters gives $\delta F=\delta G=0$.

Define normal-form coefficient derivatives by
\[
 \theta_u(u^at^iz^j)=a u^{a-1}t^iz^j,\qquad
 \theta_z(u^at^iz^j)=j u^at^iz^{j-1}.
\]
The first expression is zero when $a=0$.
These maps have degrees $2,-2$ and satisfy
\begin{equation}\label{v4-arith:ks:derivative-defects}
 \delta\theta_u=-2zF,\qquad
 \delta\theta_z=-2uF+2zG.
\end{equation}
For a direct verification take $x=u^at^iz^j$ and $y=u^bt^\ell z^m$.
Put $A=a+b$, $I=i+\ell$, and $J=j+m$.
The formulas are
\[
 F(x,y)=
 \begin{cases}
 u^At^{I-r}z^J,&I\geq r,\ J\leq1,\\
 0,&\text{otherwise},
 \end{cases}
\]
\[
 G(x,y)=
 \begin{cases}
 u^At^I,&I<r,\ J=2,\\
 2u^{A+1}t^{I-r}z^{J-1},&I\geq r,\ J\geq1,\\
 0,&\text{otherwise}.
 \end{cases}
\]
Substitution proves \eqref{v4-arith:ks:derivative-defects} over $S$.
No derivative of $t^r$ occurs, so no inverse of $r$ is required.

For even internal degrees use the cup product
$(P\smile Q)(\mathbf x,\mathbf y)=P(\mathbf x)Q(\mathbf y)$.
For the two binary cochains put
\[
 (F\circ G)(x,y,z)=F(G(x,y),z)-F(x,G(y,z)).
\]
Expansion of the Hochschild differential gives
\[
 \delta(F\circ G)=G\smile F-F\smile G.
\]
It follows that the ternary cochain of degree $-2$
\begin{equation}\label{v4-arith:ks:fourth-primitive}
 g=\tfrac12(\theta_z\smile F+\theta_u\smile G)-z(F\circ G)
\end{equation}
satisfies $\delta g=-uF\smile F=m_4$.
The two terms with coefficient $z$ cancel by the preceding identity.

Let $Q$ be the tensor-coalgebra coordinate change with components
$q_1=1$, $q_3=-g$, and all other $q_n=0$.
It is invertible because its deviation from the identity lowers tensor length.
The inverse is a finite sum on each finite tensor word.
Conjugating the transferred coderivation gives a minimal structure $m'$ with
\[
 m'_3=m'_4=m'_5=0,
\]
\[
 m'_6=m_6-\sum_{j=0}^3
 m_4(1^{\otimes j}\otimes g\otimes1^{\otimes(3-j)})+g\smile g.
\]
The arity-four identity uses $m'_4=m_4-\delta g$.
The arity-six signs are the morphism signs specified in
Section~\ref{v4-arith:ks:fourth-operation}.
Applying the augmentation gives the decisive identity
\begin{equation}\label{v4-arith:ks:sixth-augmentation}
 \epsilon m'_6=(\epsilon g)\smile(\epsilon g).
\end{equation}
The arity-seven identity makes $m'_6$ a Hochschild cocycle.

\begin{lemma}[a bar cycle detects the sixth operation]
\label{v4-arith:ks:sixth-pairing}
\ClaimStatusProvedHere{}
For $r\geq2$, the cocycle $m'_6$ is not a Hochschild boundary.
\end{lemma}

\begin{proof}
Put $a=t^{r-1}$.
In the reduced bar complex of the augmented algebra $H_C$, use
\[
 b[x_1|\cdots|x_n]
 =\sum_{j=1}^{n-1}(-1)^{j-1}
 [x_1|\cdots|x_jx_{j+1}|\cdots|x_n].
\]
All letters have even internal degree, so their suspensions are odd.
The chains
\[
 \rho=[a|t]-[u|z]-[z|u],\qquad \sigma=[z|z]
\]
are cycles: $b\rho=t^r-2uz=0$ and $b\sigma=z^2=0$.
Write $\mathrm{sh}$ for the signed shuffle product.
It interleaves words while preserving each word's order.
Its sign is the sign of the permutation of the suspended letters.
It commutes with $b$ because products between different words cancel in pairs,
while products inside one word give its own bar differential.
Thus $\gamma=\rho\mathbin{\mathrm{sh}}\rho\mathbin{\mathrm{sh}}\sigma$ is a length-six bar cycle.

On its letters, \eqref{v4-arith:ks:fourth-primitive} gives
\[
 \epsilon g(z,a,t)=\epsilon g(z,t,a)=\tfrac12,
 \qquad \epsilon g(u,z,z)=\tfrac12.
\]
All other triples contribute zero.
For $r=2$ the equality $a=t$ identifies the first two triples.
The complete list of nonzero contributions to
$((\epsilon g)\smile(\epsilon g))(\gamma)$ is
\[
 \begin{array}{c|r|r}
 \text{word}&\text{shuffle coefficient}&\text{cochain value}\\ \hline
 \bigl[z|a|t|z|a|t\bigr]&2&1/4\\
 \bigl[z|a|t|u|z|z\bigr]&-2&1/4\\
 \bigl[u|z|z|z|a|t\bigr]&-2&1/4
 \end{array}
\]
These coefficients are obtained by choosing the two ordered positions
of each of the three two-letter words and summing their permutation signs.
The table remains valid after $a=t$, with identical words collected.
Therefore
\[
 \epsilon m'_6(\gamma)=\tfrac12-\tfrac12-\tfrac12=-\tfrac12.
\]
If $m'_6=\delta K$, its exterior Hochschild terms vanish after $\epsilon$
because every bar letter has zero augmentation.
Its interior terms evaluate $K$ on $b\gamma=0$, up to an overall sign.
Their sum is zero.
The element $-1/2$ is a unit in the nonzero ring $S$, giving the contradiction.
\end{proof}

\begin{theorem}[formality of the full relative centre]
\label{v4-arith:ks:full-centre-formality}
\ClaimStatusProvedHere{}
For every commutative unital ring $S$ and every $r\geq1$,
\[
 C_{S,r}\text{ is formal over }S
 \quad\Longleftrightarrow\quad r=1\ \text{or}\ S=0.
\]
For a nonzero $S$ and $r\geq2$, it also has no formality zigzag
that fixes $S[t]$.
\end{theorem}

\begin{proof}
For $r=1$, the inclusion $S[u,z]/(z^2)\hookrightarrow C_{S,1}$
is a quasi-isomorphism by monic division.
For $S=0$ the statement is immediate.
Suppose $r\geq2$ and $S\ne0$.
If $2$ is a nonunit, a formality zigzag would truncate to one for
$A(S,2,r)$, contradicting Theorem~\ref{v4-arith:ks:coefficient-formality}.

If $2$ is a unit, use the minimal structure $m'$ constructed above.
Every homogeneous component of $H_C$ is free over $S$.
The finite comparison lemma transports a hypothetical formality zigzag
through arity six to a morphism from $m'$ to strict $H_C$.
Normalize its linear component to the identity.
Parity forces its second and fourth components to vanish.
Let $a_3,a_5$ denote its third and fifth components.
The arity-four and arity-six identities give
\[
 \delta a_3=0,\qquad m'_6=\delta a_5+a_3\smile a_3.
\]
The self-insertion
\[
 \begin{split}
 (a_3\circ a_3)(x_1,\ldots,x_5)={}&a_3(a_3(x_1,x_2,x_3),x_4,x_5)\\
 &+a_3(x_1,a_3(x_2,x_3,x_4),x_5)\\
 &+a_3(x_1,x_2,a_3(x_3,x_4,x_5))
 \end{split}
\]
satisfies $\delta(a_3\circ a_3)=-2a_3\smile a_3$.
To verify this identity, expand each differential by the displayed Hochschild rule.
The terms with an adjacent multiplication inside an inner or outer $a_3$
cancel by $\delta a_3=0$.
The two remaining end decompositions each give $-a_3\smile a_3$.
Thus $m'_6=\delta(a_5-\tfrac12a_3\circ a_3)$,
contrary to Lemma~\ref{v4-arith:ks:sixth-pairing}.
Finally, forgetting the extra $S[t]$-action would turn any relative
formality zigzag into an $S$-linear one.
\end{proof}

The full centre therefore remains nonformal over every field when $r\geq2$.
At characteristic two its fourth operation is obstructed.
When two is invertible, the fourth operation is exact and the sixth operation is obstructed.
The nonpositive truncation has the different coefficient criterion of
Theorem~\ref{v4-arith:ks:coefficient-formality}.

\section{Coefficient change and the arithmetic comparison}
\label{v4-arith:ks:coefficient-change}

For a ring map $S\to S'$, the normal bases and monic-division contraction give
\[
 A(S,c,r)\otimes_S S'=A(S',c',r),\qquad
 H^*(A(S,c,r))\otimes_S S'\cong H^*(A(S',c',r)),
\]
where $c'$ is the image of $c$.
The same argument gives
\[
 C_{S,r}\otimes_S S'=C_{S',r},\qquad
 H^*(C_{S,r})\otimes_S S'\cong H^*(C_{S',r}),
\]
with
\[
 H^*(C_{S,r})=S[u,t,z]/(z^2,t^r-2uz).
\]
These are assertions for the displayed complexes and their split normal forms.
They hold without flatness of $S'$.
The relative-centre resolution also extends to $S'[t]$ on its two free generators,
so its multiplicative identification persists after coefficient change.

\begin{corollary}[the exceptional coefficient]
\label{v4-arith:ks:exceptional-two}
\ClaimStatusProvedHere{}
For $r\geq2$, the nonpositive part of the integral relative centre is nonformal.
Its extension from $\mathbb Z$ to $S$ is formal exactly when $2\in S^\times$.
In particular it is formal over $\mathbb Q$ and every field of odd characteristic,
and nonformal over every field of characteristic two and over $\mathbb Z/4$.
The full $C_{S,r}$ is nonformal whenever $2$ is a nonunit in $S$ and $r\geq2$.
For $r=1$, both the full centre and its nonpositive part are formal over every $S$.
\end{corollary}

\begin{proof}
The assertions about the nonpositive part follow from
Proposition~\ref{v4-arith:ks:centre-truncation} and Theorem~\ref{v4-arith:ks:coefficient-formality}.
For any cochain DGA $D$, the complex with $D^n$ in negative degrees,
$Z^0(D)$ in degree zero, and zero in positive degrees is a sub-DGA.
This truncation functor preserves quasi-isomorphisms.
Applying it to a formality zigzag for $C_{S,r}$ would give a formality
zigzag for $A(S,2,r)$.
This proves the stated nonformality of the full centre.
For $r=1$, the inclusion
$S[u,z]/(z^2)\hookrightarrow C_{S,1}$ is a quasi-isomorphism:
monic division identifies cohomology by substituting $t=2uz$.
The analogous inclusion of $B$ proves the nonpositive assertion.
\end{proof}

\begin{corollary}[localization of the full centre]
\label{v4-arith:ks:centre-localization}
\ClaimStatusProvedHere{}
If $2\in S^\times$, then $C_{S,r}[u^{-1}]$ is formal over
$S[t,u,u^{-1}]$ for every $r\geq1$.
\end{corollary}

\begin{proof}
After localization, $z=u^{-1}v$ gives the inverse of
\[
 A(S,2,r)[u^{-1}]\xrightarrow{\ j\ }C_{S,r}[u^{-1}].
\]
The zigzag in Proposition~\ref{v4-arith:ks:positive-zigzag} fixes $S[t,u]$.
Localization at its central even element $u$ is exact on modules,
so it preserves both quasi-isomorphisms.
Its cohomology target is
$S[t,u,u^{-1}]/(t^{2r})$, with $z=t^r/(2u)$.
\end{proof}

\begin{example}[unit and nonunit parameters]
For $S=\mathbb F_2$ and $r=2$, the model $A(S,0,2)$ has
$m_4(t,t,t,t)=-u=u$ and is nonformal.
The model $A(S,1,2)$ is formal by the explicit zigzag.
It is not the reduction of $A(\mathbb Z,2,2)$, whose parameter reduces to zero.
Over $\mathbb Z/4$, the parameters $2$ and $3$ likewise give
nonformal and formal models, respectively.
The criterion is invertibility of $c$, rather than its nonvanishing.
For the zero ring, the unique element is a unit and every displayed algebra is zero.
\end{example}

The global arithmetic object
$\mathcal V^{\mathrm{prim}}$
of Definition~\ref{v4-arith:def:global-iwahori-hecke} is a Hochschild trace class.
The centre above is a relative derived endomorphism algebra.
A comparison between them requires a multiplicative map with a specified base,
source, target, and cohomology isomorphism.
An equality of degree-zero centres does not supply these data.

\begin{proposition}[scope of the centre obstruction]
\label{v4-arith:ks:centre-comparison-scope}
\ClaimStatusProvedHereConditional{}
Let $\mathcal Z$ be an associative $S$-DGA model for an arithmetic derived centre.
Suppose a finite multiplicative $S$-linear quasi-isomorphism zigzag identifies
$\tau_{\leq0}\mathcal Z$ with $A(S,2,r)$.
If $r\geq2$ and $2$ is a nonunit, then $\mathcal Z$ is nonformal over $S$.
If the full $\mathcal Z$ is multiplicatively quasi-isomorphic to $C_{S,r}$,
it is nonformal for every nonzero $S$ and $r\geq2$.
A relative dg Morita equivalence with $E_r$ over $S[t]$, together with
its bimodule comparison, is one possible source of the required identification.
\end{proposition}

\begin{proof}
A formality zigzag for $\mathcal Z$ would truncate to one for
$\tau_{\leq0}\mathcal Z$.
Composing with the hypothesized comparison would make $A(S,2,r)$ formal,
contrary to Theorem~\ref{v4-arith:ks:coefficient-formality}.
The full comparison transfers the classification in
Theorem~\ref{v4-arith:ks:full-centre-formality}.
\end{proof}

The positive zigzag fixes $S[t,u]$ on the nonpositive algebra.
The full centre has the additional sixth-operation obstruction of
Theorem~\ref{v4-arith:ks:full-centre-formality}.
For the global Iwahori category, the required relative bimodule comparison
is an additional construction, independent of rectification of the $E_1$ operad.


\section{The Chiral Bracket via Adelic Residue}
\label{v4-arith:ks:chiral-bracket-adelic}

\begin{definition}[archimedean \(\lambda\)-bracket on the cover]
\label{v4-arith:def:arith-lambda-bracket}
Let \(\overline{C}\) denote the Arakelov compactification of the
arithmetic curve, obtained by adjoining the archimedean place
\(|\cdot|_{\infty}\) to \(\mathrm{Spec}\,\mathbb{Z}\).
The Beilinson--Drinfeld \(\lambda\)-bracket is defined only on the
archimedean Heisenberg cover, where a holomorphic diagonal and a
complex residue exist:
\begin{equation}
\{a,b\}_{\infty}\;=\;\mathrm{Res}_{z=w}\,\omega_{a,b}(z,w).
\end{equation}
No finite prime contributes a vertex-algebra residue on
\(\mathcal{V}^{\mathrm{prim}}\). A finite-place bracket would require
a separate prismatic construction and is not part of the Hochschild
trace class.
\end{definition}

\begin{theorem}[finite-place residue obstruction]
\label{v4-arith:thm:chiral-bracket-adelic}
\ClaimStatusProvedHere{} The archimedean bracket above does not descend
to a Beilinson--Drinfeld bracket on
\(\mathcal{V}^{\mathrm{prim}}\) over \(\overline{C}\). Finite places are
totally disconnected in the factorization topology used here, and the
global object is an \(E_{1}\) Hochschild trace class rather than a
vertex algebra.
\end{theorem}

\section{Euler logarithms and prime traces}
\label{v4-arith:ks:logL-not-cyclic}

An Euler product already has a concrete algebraic carrier: the free
commutative monoid on the primes. Its logarithm singles out prime powers,
and an energy derivation supplies the factors $\log p$. These two
operations must be distinguished before seeking a Hall-algebra realization.

Let $\mathscr D_K$ be the algebra of formal Dirichlet series over a
commutative $\mathbb Q$-algebra $K$. An element is a family
$a=\sum_{n\geq1}a_n[n]$, with multiplication $[m][n]=[mn]$;
the coefficient of $[n]$ in $ab$ is $\sum_{d\mid n}a_db_{n/d}$.
We use the coefficient topology, or equivalently the filtration by the
ideals of series supported on $n\geq N$. If $a_1=0$, then the coefficient
of $[n]$ in $a^j$ vanishes for $2^j>n$, so formal logarithms and
exponentials converge in this topology.

\begin{proposition}[Euler logarithm and energy derivation]
\label{v4-arith:obs:logL-not-cyclic}
\ClaimStatusProvedHere{}
In $\mathscr D_{\mathbb Q}$ one has
\begin{equation}
 Z=\prod_p(1-[p])^{-1}=\sum_{n\geq1}[n],
 \qquad
 \log Z=\sum_p\sum_{k\geq1}\frac{[p^k]}{k}.
 \label{v4-arith:ks:euler-log}
\end{equation}
After extending coefficients to $\mathbb R$, the continuous derivation
$D[n]=(\log n)[n]$ satisfies
\begin{equation}
 D\log Z=\sum_{n\geq1}\Lambda(n)[n],\qquad
 \Lambda(n)=
 \begin{cases}\log p&n=p^k,\ k\geq1,\\0&\text{otherwise.}\end{cases}
 \label{v4-arith:ks:euler-derivation}
\end{equation}
For $\operatorname{Re}s>1$, their convergent analytic realizations are
\begin{equation}
 \log\zeta(s)=\sum_{p,k\geq1}\frac{p^{-ks}}{k},
 \qquad
 -\frac{\zeta'(s)}{\zeta(s)}
   =\sum_{p,k\geq1}(\log p)p^{-ks}.
 \label{v4-arith:ks:euler-analytic}
\end{equation}
The logarithm is the holomorphic branch tending to zero as
$\operatorname{Re}s\longrightarrow+\infty$.
\end{proposition}

\begin{proof}
Unique prime factorization gives exactly one contribution to every
coefficient of the product in \eqref{v4-arith:ks:euler-log}.
Taking its logarithm coefficient by coefficient gives
$-\log(1-[p])=\sum_{k\geq1}[p^k]/k$.
The identity $\log(mn)=\log m+\log n$ proves the Leibniz rule for $D$.
Since $D[p^k]=k\log p\,[p^k]$, applying $D$ cancels the denominator $k$.

On a compact subset of $\operatorname{Re}s>1$, choose $\sigma_0>1$
below all real parts. Then
\[
 \sum_{p,k\geq1}\frac{|p^{-ks}|}{k}
 \leq\frac{1}{1-2^{-\sigma_0}}\sum_{n\geq2}n^{-\sigma_0}<\infty.
\]
The differentiated series has the same bound with
$\sum_{n\geq2}(\log n)n^{-\sigma_0}$ in place of the final sum.
Thus both series converge locally uniformly. Exponentiation identifies
the first series with the Euler product for $\zeta$ and shows that this
product is nonzero. Differentiation proves the second identity.
\end{proof}

The coefficient of $n^{-s}$ in $\log\zeta$ is therefore $1/k$ for
$n=p^k$, and zero otherwise. It is rational. The factor $\log p$
belongs to the differentiated series. Neither assertion involves a
cyclic-cohomology construction.

\begin{proposition}[local determinant factors]
\label{v4-arith:ks:local-determinant-traces}
\ClaimStatusProvedHere{}
For each prime $p$, let $T_p$ be a complex endomorphism of a Hermitian
space of dimension at most $d$, where $d$ is fixed. Suppose
$\|T_p\|\leq p^\theta$ for a fixed $\theta\geq0$. On
$\operatorname{Re}s>1+\theta$ the product
\[
 L(s)=\prod_p\det(1-p^{-s}T_p)^{-1}
\]
is holomorphic and nonzero, and
\begin{align}
 \log L(s)&=\sum_{p,k\geq1}
      \frac{\operatorname{Tr}(T_p^k)}{k}p^{-ks},
      \label{v4-arith:ks:matrix-euler-log}\\
 -\frac{L'(s)}{L(s)}&=\sum_{p,k\geq1}
      (\log p)\operatorname{Tr}(T_p^k)p^{-ks}.
      \label{v4-arith:ks:matrix-euler-derivative}
\end{align}
Here $\operatorname{Tr}$ is the ordinary, unnormalized matrix trace,
and the logarithm tends to zero as the real part of $s$ tends to infinity.
\end{proposition}

\begin{proof}
For $|q|\|T\|<1$, Jacobi's determinant formula gives
\[
 \frac{d}{dq}\bigl[-\log\det(1-qT)\bigr]
   =\operatorname{Tr}\bigl(T(1-qT)^{-1}\bigr)
   =\sum_{k\geq1}\operatorname{Tr}(T^k)q^{k-1}.
\]
Integration from $q=0$ fixes the constant and gives the local logarithm.
The estimate $|\operatorname{Tr}(T_p^k)|\leq d p^{k\theta}$ reduces
normal convergence, including differentiation, to the preceding proof
with $\operatorname{Re}s-\theta$ in place of $\operatorname{Re}s$.
\end{proof}

These formulas give an arithmetic use for a proposed graded carrier:
its trace must produce the specified local factors. For $T_p=1$, that
carrier can be constructed without a quiver.

\begin{proposition}[the prime monoid and its thermal trace]
\label{v4-arith:ks:prime-fock-trace}
\ClaimStatusProvedHere{}
Let $\mathcal P=\mathbb Q[x_p:p\text{ prime}]$, with each monomial
involving finitely many variables. Declare its monomials orthonormal
after extension to $\mathbb C$, and complete to a Hilbert space
$\mathcal F$. The assignment
\[
 x^\nu\longmapsto\varepsilon_{\prod_p p^{\nu_p}}
\]
is a unitary identification $\mathcal F\cong\ell^2(\mathbb N_{\geq1})$.
The diagonal energy operator
\[
 Hx^\nu=\left(\sum_p\nu_p\log p\right)x^\nu
\]
is nonnegative and self-adjoint on
$\{\sum_n a_n\varepsilon_n:\sum_n(\log n)^2|a_n|^2<\infty\}$.
For $\operatorname{Re}s>1$,
\begin{equation}
 \operatorname{Tr}(e^{-sH})=\zeta(s),\qquad
 \operatorname{Tr}(He^{-sH})=-\zeta'(s).
 \label{v4-arith:ks:prime-thermal-trace}
\end{equation}
For real $\beta>1$, the normalized state
$\langle A\rangle_\beta=\operatorname{Tr}(Ae^{-\beta H})/\zeta(\beta)$,
whenever the numerator is defined, satisfies
\[
 \langle H\rangle_\beta=-\frac{\zeta'(\beta)}{\zeta(\beta)},
 \qquad
 \langle H^2\rangle_\beta-\langle H\rangle_\beta^2
   =\sum_{p,k\geq1}k(\log p)^2p^{-k\beta}.
\]
\end{proposition}

\begin{proof}
The monomial basis is indexed bijectively by positive integers through
unique factorization. Under this identification, $H\varepsilon_n
=(\log n)\varepsilon_n$. A real diagonal multiplication operator on
its maximal square-summability domain is self-adjoint; this follows
directly by testing the adjoint against each basis vector.
The singular values of $e^{-sH}$ are $n^{-\operatorname{Re}s}$, so it
is trace class on the stated half-plane. The same argument applies to
$H^je^{-sH}$ for every fixed $j\geq0$. Its diagonal trace proves
\eqref{v4-arith:ks:prime-thermal-trace} and permits two differentiations.
The quotient rule then identifies the variance with
$(\log\zeta)''(\beta)$, and \eqref{v4-arith:ks:euler-analytic} evaluates it.
\end{proof}

For a finite set of primes $F$, the same construction on
$\mathcal P_F=\mathbb Q[x_p:p\in F]$ has partition function
$\zeta_F(s)=\prod_{p\in F}(1-p^{-s})^{-1}$ for
$\operatorname{Re}s>0$. Multiplication by $x_p$ extends to the isometry
$\varepsilon_n\mapsto\varepsilon_{pn}$. Thus the monoid action gives
the prime shifts in the Bost--Connes representation. The displayed
trace is a Hilbert-space trace; a map to the arithmetic Hochschild
trace class $\mathcal V^{\mathrm{prim}}$ is additional data.

\section{Hall multiplication and the BPS filtration}
\label{v4-arith:ks:CoHA-conjectural}

A Hall product counts extensions of representations. A monoid product
multiplies their prime labels. Even for one vertex these operations
produce different algebras, and the distinction persists when their
graded characters resemble Euler products.

\begin{proposition}[the zero-loop and one-loop calculations]
\label{v4-arith:ks:one-vertex-coha}
\ClaimStatusProvedHere{}
Let $Q_m$ have one vertex and $m\in\{0,1\}$ loops, with zero potential.
Its rational CoHA has dimension-$n$ summand
$\mathcal H^{(m)}_n=\mathbb Q[x_1,\ldots,x_n]^{\mathfrak S_n}$.
A homogeneous polynomial of polynomial degree $j$ has cohomological
degree $2j+(1-m)n^2$. Write $e_k=x^k\in\mathcal H^{(m)}_1$.
Then
\[
 \mathcal H^{(0)}\cong\bigwedge_{\mathbb Q}(e_0,e_1,e_2,\ldots),
 \quad |e_k|=2k+1,
 \qquad
 \mathcal H^{(1)}\cong\mathbb Q[e_0,e_1,e_2,\ldots],
 \quad |e_k|=2k.
\]
Every generator has dimension degree one.
\end{proposition}

\begin{proof}
The representation space is affine and equivariantly contractible,
so its equivariant cohomology is $H^*(B\mathrm{GL}_n,\mathbb Q)$.
The Hall correspondence chooses an invariant subspace. Localization
at the diagonal torus gives the shuffle formula
\begin{equation}
 f*g=\sum_{\substack{I\sqcup J=\{1,\ldots,a+b\}\\|I|=a}}
 f(x_I)g(x_J)\prod_{i\in I,\,j\in J}(x_j-x_i)^{m-1}.
 \label{v4-arith:ks:one-vertex-shuffle}
\end{equation}
The numerator is the Euler class of the $m$ arrow blocks and the
denominator the Euler class of the tangent space to the Grassmannian.
This is the specialization of Kontsevich--Soibelman,
\emph{Cohomological Hall algebra, exponential Hodge structures and
motivic Donaldson--Thomas invariants}, Theorem~2, pp.~11--12.

For $m=0$, the product of $n$ dimension-one elements is their
alternant divided by $\prod_{i<j}(x_j-x_i)$. Repeated exponents give
zero. For $0\leq k_1<\cdots<k_n$, it is the Schur polynomial of
partition $(k_n-n+1,k_{n-1}-n+2,\ldots,k_1)$. These form a basis of
the symmetric polynomials. Hence the exterior-algebra map is an
isomorphism. In particular,
\[
 e_0*e_0=0,\qquad
 e_0*e_1=\frac{x_2-x_1}{x_2-x_1}=1\in\mathcal H^{(0)}_2,
 \qquad e_1*e_0=-1\in\mathcal H^{(0)}_2.
\]
The constants on the right lie in dimension degree two and are not
the unit, which lies in dimension degree zero.

For $m=1$, the shuffle factor is one. Products of the $e_k$ give the
monomial symmetric basis, multiplied by the factorials of repeated
exponents. Those factorials are invertible over $\mathbb Q$, proving
the polynomial-algebra assertion. This argument also explains its
coefficient restriction: over $\mathbb Z$, for example,
$e_0*e_0=2\cdot1\in\mathcal H^{(1)}_2$.
\end{proof}

An edgeless quiver with a vertex for each prime in a nonempty finite
set $F$ therefore has odd square-zero classes already in each
one-vertex subalgebra. Its CoHA cannot be $\mathcal P_F$.
For $F=\varnothing$, there is only the zero representation, so both
the CoHA and $\mathcal P_{\varnothing}$ are $\mathbb Q$.
For the disjoint union of Jordan quivers, the rational Hall algebra
is polynomial on generators $e_{p,k}$ for $p\in F$, $k\geq0$.
There is, however, a precise prime-monoid comparison:
\begin{equation}
 \mathcal P_F\ \xrightarrow{\ x_p\mapsto e_{p,0}\ }\
 \mathcal H_F^{(1)},\qquad
 \mathcal H_F^{(1)}/(e_{p,k}:k>0)\cong\mathcal P_F.
 \label{v4-arith:ks:jordan-prime-maps}
\end{equation}
The first map identifies the cohomological-degree-zero subalgebra;
the second kills the positive-degree generators. The maps are
compatible with inclusions of finite sets of primes. The full local
graded character, as a series in $\mathbb Q[[z,q]]$ with $q$ recording
half the cohomological degree, is
\[
 \sum_{n,j\geq0}\dim\bigl(\mathcal H^{(1)}_{n,2j}\bigr)z^nq^j
   =\prod_{k\geq0}(1-zq^k)^{-1}.
\]
Its coefficient of $q^0$ is $(1-z)^{-1}$, the prime Euler factor
after $z=p^{-s}$. Forgetting the $q$-grading instead gives infinite
multiplicities even in dimension degree one. Thus
\eqref{v4-arith:ks:jordan-prime-maps} specifies the sector on which
Proposition~\ref{v4-arith:ks:prime-fock-trace} applies.

To include a potential, fix a finite quiver $Q$ with symmetric arrow
matrix and a polynomial cyclic potential
$W\in\mathbb C Q/[\mathbb C Q,\mathbb C Q]_{\mathrm{vect}}$.
For a dimension vector $\mathbf d\in\mathbb N^{Q_0}$, put
\[
 X_{\mathbf d}=\bigoplus_{a:i\to j}
     \operatorname{Hom}(\mathbb C^{d_i},\mathbb C^{d_j}),\qquad
 G_{\mathbf d}=\prod_i\mathrm{GL}_{d_i}(\mathbb C),\qquad
 \mathfrak M_{\mathbf d}=[X_{\mathbf d}/G_{\mathbf d}].
\]
The Euler form is
$\chi(\mathbf d,\mathbf e)=\sum_i d_ie_i-\sum_{a:i\to j}d_ie_j$,
so $\dim\mathfrak M_{\mathbf d}=-\chi(\mathbf d,\mathbf d)$.
Evaluation followed by matrix trace defines the invariant function
$f_{\mathbf d}=\operatorname{Tr}(W)$ on this stack and on its affine
coarse moduli space $M_{\mathbf d}=X_{\mathbf d}/\!/G_{\mathbf d}$.
Assume $\operatorname{crit}(f_{\mathbf d})\subset f_{\mathbf d}^{-1}(0)$
for every $\mathbf d$; zero and homogeneous positive-degree potentials
satisfy this condition.

We use rational monodromic mixed Hodge structures and the monodromic
vanishing-cycle functor $\phi_f^{\mathrm{mon}}$, normalized to be
perverse exact and to satisfy $\phi_0^{\mathrm{mon}}=\mathrm{id}$.
Let $\mathbb L=H_c^*(\mathbb A^1,\mathbb Q)$, in cohomological
degree two. Its monodromic square root is
$\mathbb L^{1/2}=H_c^*(\mathbb A^1,\phi_{x^2}^{\mathrm{mon}}
\mathbb Q_{\mathbb A^1})$, in degree one, with the square-root
isomorphism fixed by the complex orientation. On a smooth space or smooth quotient
stack $Y$, write
$\mathrm{IC}_Y^{\mathrm{vir}}=\mathbb L^{-\dim Y/2}\otimes\mathbb Q_Y$.
On a singular irreducible coarse moduli space, use the intersection
complex with the same normalization on its smooth locus and
intermediate extension. These conventions include the shifts in the
one-vertex calculation.

\begin{definition}[critical cohomological Hall algebra]
\label{v4-arith:thm:KS-CoHA}
The underlying dimension-graded, cohomologically graded object is
\begin{equation}
 \mathcal H_{Q,W}
 =\bigoplus_{\mathbf d}
 \left[
 H_c^*(\mathfrak M_{\mathbf d},
       \phi_{f_{\mathbf d}}^{\mathrm{mon}}
       \mathrm{IC}_{\mathfrak M_{\mathbf d}}^{\mathrm{vir}})
 \right]^\vee.
 \label{v4-arith:ks:normalized-coha}
\end{equation}
Here $\vee$ denotes graded duality, including duality of Hodge
structures, and equivariant cohomology is defined by finite-dimensional
approximations to the classifying spaces. Its associative product is
the vanishing-cycle pullback and proper pushforward along the stack
of short exact sequences, with the virtual shifts displayed above.
The equality of trace potentials on a block upper-triangular
representation and on its subquotients supplies the required
Thom--Sebastiani map. The zero representation supplies the unit.
This is the critical CoHA of Kontsevich--Soibelman, \S7, in the
normalization of Davison--Meinhardt, equation~(1) and \S5.2.
\end{definition}

Choose a stability condition with all central charges $\zeta_i=\sqrt{-1}$.
Every representation is semistable of slope zero, since the slope
of $\mathbf d\ne0$ is
$-\operatorname{Re}(\sum_i\zeta_i d_i)/
\operatorname{Im}(\sum_i\zeta_i d_i)=0$.
This stability is slope-generic: the antisymmetrization of $\chi$
vanishes on all dimension vectors because $Q$ is symmetric.
Thus the single-slope integrality theorem applies to the full algebra
\eqref{v4-arith:ks:normalized-coha}.

For $\mathbf d\ne0$, define a BPS sheaf on $M_{\mathbf d}$ by
\[
 \mathcal B_{\mathbf d}=
 \begin{cases}
 \phi_{f_{\mathbf d}}^{\mathrm{mon}}
       \mathrm{IC}_{M_{\mathbf d}}^{\mathrm{vir}}
       &\text{if the stable locus is nonempty},\\
 0&\text{otherwise},
 \end{cases}
 \qquad
 B=\bigoplus_{\mathbf d\ne0}H^*(M_{\mathbf d},\mathcal B_{\mathbf d}).
\]
At this stability, stable representations are precisely the simple
ones. Finally put
\[
 U=H^*(B\mathbb C^*,\mathbb Q)_{\mathrm{vir}}\otimes B
   =\mathbb L^{1/2}\otimes\mathbb Q[u]\otimes B,
 \qquad |u|=2.
\]
Here $u$ is the first Chern class of the universal line bundle and
has dimension degree zero. In particular, the BPS
cohomology $B$, the shifted BPS Lie algebra $\mathbb L^{1/2}\otimes B$,
and the full generator object $U$ are distinct.

\begin{theorem}[Davison--Meinhardt integrality and PBW]
\label{v4-arith:thm:davison-integrality}
\ClaimStatusProvedElsewhere{}
For the preceding data define, modulo two,
\[
 \tau(\mathbf d,\mathbf e)=\chi(\mathbf d,\mathbf e)
       +\chi(\mathbf d,\mathbf d)\chi(\mathbf e,\mathbf e).
\]
Choose a bilinear form $\psi$ with
$\psi(\mathbf d,\mathbf e)+\psi(\mathbf e,\mathbf d)
=\tau(\mathbf d,\mathbf e)$, and twist the Hall product by
$a*_\psi b=(-1)^{\psi(\mathbf d,\mathbf e)}a*b$.
Write $\mathcal H^\psi$ for the resulting algebra. It has an
exhaustive multiplicative ambient perverse filtration $P$ and a
canonical generator embedding $\iota:U\hookrightarrow\mathcal H^\psi$
such that the symmetrized Hall-product map
\begin{equation}
 \operatorname{PBW}:\operatorname{Sym}_{\mathrm{sup}}(U)
       \xrightarrow{\ \sim\ }\mathcal H^\psi
 \label{v4-arith:ks:pbw-object-map}
\end{equation}
is an isomorphism of dimension-graded, cohomologically graded
monodromic mixed Hodge structures. The associated graded
$\operatorname{gr}_P\mathcal H^\psi$ is a free supercommutative
algebra. The positive-dimension part of $P_1\mathcal H^\psi$ is
$\mathbb L^{1/2}\otimes B$, and is closed under the Hall
supercommutator.
\end{theorem}

Here $\operatorname{Sym}_{\mathrm{sup}}$ is polynomial on the even
part and exterior on the odd part, with the cohomological Koszul sign.
The form $\tau$ is alternating modulo two, so a form $\psi$ exists
by choosing an ordering of a basis of $(\mathbb Z/2)^{Q_0}$.
Bilinearity makes the twisted product associative. The theorem is
Davison--Meinhardt, \emph{Cohomological Donaldson--Thomas theory of
a quiver with potential and quantum enveloping algebras},
Theorems~A and~C, pp.~5--6 and 10--11; the Lie subalgebra is
\S1.7, equation~(22).

For clarity, $P$ uses the semisimplification map
$p_{\mathbf d}:\mathfrak M_{\mathbf d}\to M_{\mathbf d}$.
One approximates this map by the proper maps from moduli spaces
of sufficiently framed representations to $M_{\mathbf d}$, with
their virtual shifts. Perverse truncation of these direct images
on the ambient space $M_{\mathbf d}$, followed by cohomology and
stabilization in each degree, defines $P$; equivalently one uses
the dual compact-support construction with truncation
$\tau_{\geq-p}$ in Davison--Meinhardt, \S5.3, pp.~48--49.
The summand $\mathbb L^{1/2}u^k\otimes B$ supplies generators of
perverse degree $2k+1$. Supercommutativity of the associated graded
implies $[P_a,P_b]\subset P_{a+b-1}$, and hence the assertion about
$P_1$. No integral lattice or coefficient-change theorem is asserted
by this rational Hodge-theoretic construction.

\begin{proposition}[the multiplicative PBW criterion]
\label{v4-arith:ks:pbw-multiplicative-criterion}
\ClaimStatusProvedHere{}
The map \eqref{v4-arith:ks:pbw-object-map} is an algebra isomorphism
if and only if the embedded generators $\iota(U)$ pairwise
supercommute in $\mathcal H^\psi$. An ordinary polynomial algebra
on these homogeneous generators additionally requires $U$ to have
no odd part.
\end{proposition}

\begin{proof}
If the generators supercommute, the universal property gives an
algebra map from $\operatorname{Sym}_{\mathrm{sup}}(U)$ to
$\mathcal H^\psi$. On a product of $n$ generators it equals the
signed average of the $n!$ ordered Hall products, hence equals the
PBW map and is bijective. Conversely, an algebra map extending
$\iota$ preserves the defining supercommutation relations.
Over $\mathbb Q$ a nonzero odd generator in the symmetric
superalgebra has square zero, whereas an ordinary polynomial
algebra over $\mathbb Q$ is reduced.
\end{proof}

The commutator condition can fail even for a symmetric quiver with
a homogeneous cubic potential.

\begin{example}[a noncommutative symmetric critical CoHA]
\label{v4-arith:ks:tripled-a2}
Start with $Q=(1\xrightarrow{a}2)$, add the reverse arrow $a^*$
and loops $\omega_1,\omega_2$, and call the resulting symmetric
quiver $\widetilde Q$. With path multiplication given by composition
from right to left, take the cyclic potential
\[
 \widetilde W=(aa^*-a^*a)(\omega_1+\omega_2).
\]
In dimension $(1,1)$ its trace is
$(\omega_2-\omega_1)aa^*$. For the sign twist choose the original
Euler form
$\psi(\mathbf d,\mathbf e)
=d_1e_1+d_2e_2-d_1e_2\pmod2$.
It obeys the required relation because
$\chi_{\widetilde Q}(\mathbf d,\mathbf e)
=-d_1e_2-d_2e_1$.

Dimensional reduction identifies the resulting critical CoHA with
the preprojective CoHA of $Q$, with this multiplication and twist.
The degree-zero BPS Lie algebra is
$\mathfrak n^-_{\mathfrak{sl}_3}$: this is Davison,
\emph{BPS Lie algebras and the less perverse filtration on the
preprojective CoHA}, \S3.2, equation~(35), and Theorem~6.6,
pp.~33--34. That theorem concerns the Lie algebra embedded in the
actual Hall algebra, not just an associated graded algebra.
Under its simple-root normalization, the two vertex classes
$\alpha_1,\alpha_2$ correspond to $E_{21},E_{32}$. Consequently
\[
 [\alpha_2,\alpha_1]\ne0,
 \qquad [E_{32},E_{21}]=E_{31}\ne0.
\]
Both classes have cohomological degree zero, so this is an ordinary
commutator of even elements. The critical CoHA is not
supercommutative, although Theorem~\ref{v4-arith:thm:davison-integrality}
gives its PBW decomposition. Thus that theorem cannot imply a
polynomial-algebra assertion for an arbitrary arithmetic potential.
\end{example}

\section{Cyclic potentials and the arithmetic comparison}
\label{v4-arith:ks:arithmetic-potential-comparison}

The logarithmic Euler series does have a cyclic realization, but its
critical locus tests whether it can serve as a Hall potential.
For an associative algebra $A$, degree-zero cyclic \emph{homology}
is $\mathrm{HC}_0(A)=A/[A,A]_{\mathrm{vect}}$.
Over a field, degree-zero cyclic cohomology consists of linear
traces $\lambda:A\to K$ satisfying $\lambda(ab)=\lambda(ba)$.
These constructions neither adjoin coefficients nor supply a map
from an arbitrary Dirichlet series into a path algebra.

\begin{proposition}[the completed logarithmic potential]
\label{v4-arith:ks:logarithmic-potential-obstruction}
\ClaimStatusProvedHere{}
For a finite set of primes $F$, let $Q_F$ be the disjoint union of
one-loop quivers, with loop $x_p$ at vertex $p$. Complete its complex
path algebra in the arrow ideal. In the quotient by the closure of
the commutator subspace, the cyclic series
\[
 W_F=\sum_{p\in F}\sum_{k\geq1}\frac{x_p^k}{k}
\]
is well defined. On tuples of finite matrices with $\|X_p\|<1$,
\[
 \operatorname{Tr}(W_F)(X)
   =-\sum_{p\in F}\log\det(1-X_p).
\]
In particular, at dimension one at each vertex and
$X_p=p^{-s}$ it equals $\log\zeta_F(s)$ for
$\operatorname{Re}s>0$. Its completed Jacobi algebra is zero.
\end{proposition}

\begin{proof}
There are no paths between different vertices, so the completed
path algebra is $\prod_{p\in F}\mathbb C[[x_p]]$.
Convergence and the trace identity follow from the local determinant
calculation. The cyclic derivative removes one occurrence of the
specified arrow and sums over its positions. Hence
\[
 \frac{\partial W_F}{\partial x_p}
   =\sum_{k\geq1}x_p^{k-1}=(1-x_p)^{-1}.
\]
This element is a unit in the corner with identity $e_p$.
The closed two-sided ideal generated by the cyclic derivatives
therefore contains every $e_p$ and hence the identity $\sum_{p\in F}e_p$.
Its quotient, the completed Jacobi algebra, is zero.
Equivalently, in every nonzero dimension vector the trace function
has no critical point on the matrix domain: its derivative in a variation $\delta X_p$ is
$\operatorname{Tr}((1-X_p)^{-1}\delta X_p)$, a nonzero covector.
\end{proof}

This series lies in a completed cyclic path algebra, whereas
Theorem~\ref{v4-arith:thm:davison-integrality} is stated for a
polynomial potential. The obstruction is already visible for every
polynomial truncation containing its linear term: in the formal
neighborhood of the zero representation, its cyclic derivative is
a unit and its vanishing cycles are zero in positive dimension.
Thus the displayed logarithmic realization supplies no nonzero
local BPS object at the origin.

Removing the low-degree terms changes the critical problem in a
computable way. At one vertex, for $m\geq2$, put
\[
 W^{\geq m}=\sum_{k\geq m}\frac{x^k}{k},\qquad
 \frac{\partial W^{\geq m}}{\partial x}
   =\frac{x^{m-1}}{1-x}.
\]
Its completed Jacobi algebra is
$\mathbb C[[x]]/(x^{m-1})$: it is $\mathbb C$ for $m=2$ and
the dual numbers for $m=3$. The scalar trace at $x=p^{-s}$ is now
$-\log(1-p^{-s})-\sum_{1\leq k<m}p^{-ks}/k$.
The missing terms must therefore be supplied by a separate
construction if this singularity is used in an Euler-factor model.

\begin{proposition}[arithmetic comparison criterion]
\label{v4-arith:thm:arithmetic-CoHA-conjectural}
\ClaimStatusConditional{}
Suppose that for every finite set of primes $F$ one has a finite
symmetric complex quiver with polynomial potential $(Q_F,W_F)$
satisfying the hypotheses above, and a dimension-graded rational
object $V_F$ with a specified isomorphism
\[
 \eta_F:V_F\xrightarrow{\ \sim\ }B_{Q_F,W_F}
\]
of cohomologically graded monodromic mixed Hodge structures.
Then PBW gives an object isomorphism
\[
 \operatorname{Sym}_{\mathrm{sup}}
    (\mathbb L^{1/2}\otimes\mathbb Q[u]\otimes V_F)
       \xrightarrow{\ \sim\ }\mathcal H_{Q_F,W_F}^{\psi}.
\]
It is an algebra isomorphism precisely when the corresponding
embedded generator tower supercommutes. If these generators
supercommute and are all even, it is an ordinary polynomial-algebra
isomorphism.
\end{proposition}

\begin{proof}
Apply Theorem~\ref{v4-arith:thm:davison-integrality} through
$\mathbb L^{1/2}\otimes\mathbb Q[u]\otimes\eta_F$ and then apply
Proposition~\ref{v4-arith:ks:pbw-multiplicative-criterion}.
\end{proof}

An arithmetic identification with $\mathcal V^{\mathrm{prim}}$
requires more than the objects $V_F$. It requires a realization of
that Hochschild trace class in the same graded coefficient category,
a comparison with the $V_F$, and maps compatible with enlarging $F$.
To obtain an Euler product, one must also specify an energy or
graded-character operation and prove that it intertwines the
comparison with the local traces of
Proposition~\ref{v4-arith:ks:local-determinant-traces}.
The polynomial algebra in that comparison contains the entire
$\mathbb Q[u]$ tower. Restricting to a single generator per prime
needs a specified subalgebra or quotient, as in
\eqref{v4-arith:ks:jordan-prime-maps}.

The construction of such quiver data and compatible comparison maps
for the arithmetic Hochschild trace remains open. The prime Fock
trace proves the Euler product on an explicit carrier; the
zero-loop calculation excludes an edgeless CoHA on a nonempty prime
set as that carrier;
the logarithmic-potential calculation excludes its most direct
completed critical realization; and the tripled-$A_2$ calculation
shows why a PBW theorem alone cannot remove the multiplication
obstruction.


\section{Arithmetic Formality Stratification}
\label{v4-arith:ks:stratification}

The operadic statements concern the chains on little intervals and little disks.
The little-interval operad is formal over $\mathbb Z$ by
Theorem~\ref{v4-arith:thm:E1-formality-automatic}.
The little-disk operad has the integral arity-two obstruction of
Theorem~\ref{v4-arith:thm:E2-fails-integrally}.
For an individual associative DGA, the multiplication on cohomology has
its own obstruction classes.
The explicit coefficient criterion is
Theorem~\ref{v4-arith:ks:coefficient-formality}.
Its centre interpretation uses the relative bimodule model and truncation map
of Proposition~\ref{v4-arith:ks:centre-model} and Proposition~\ref{v4-arith:ks:centre-truncation}.
An application to the global arithmetic centre requires the multiplicative
comparison specified in Proposition~\ref{v4-arith:ks:centre-comparison-scope}.
The archimedean residue bracket uses the separate geometric data in
Definition~\ref{v4-arith:def:arith-lambda-bracket}.
The Euler product is realized by the prime-monoid Hilbert space in
Proposition~\ref{v4-arith:ks:prime-fock-trace} and by the degree-zero
Jordan-CoHA sector in \eqref{v4-arith:ks:jordan-prime-maps}.
An identification with BPS cohomology requires the graded realization
and compatible maps of Proposition~\ref{v4-arith:thm:arithmetic-CoHA-conjectural}.
The perverse PBW theorem supplies the underlying graded object;
its multiplicative refinement requires vanishing Hall supercommutators.
